% sec_transport.tex -- Supplemental Material section for letter_v2.tex: thermal sum rules, the universal disc, the
% spectral-to-gap theorem, shear and Hall viscosity of the AdS4 brane, the higher-derivative remainder, the corrected
% viscosity bound, and what an EFT-free derivation of eta/s requires.
\section{Thermal sum rules, Hall viscosity and transport}
\label{sec:transport}

This section supports the transport and spectral statements of the Letter. Its inputs from the rest of this Supplement are
the two bounds on $\tau=(t_4^2+\alpha_T^2)^{1/2}$: $\tau\le\tau_*(\Delta_g)$ at every gap (Sec.~\ref{sec:angular}) and
$\tau\le2.8914\times10^{-8}$ at $\Delta_g\ge1000$ (Secs.~\ref{sec:full} and \ref{sec:fullpos}), each with its own inputs, and,
for one remark, the analytic all-gap bound of Sec.~\ref{sec:fpanalytic}. Grades are those of Sec.~\ref{sec:grades}. Every
computation cited here has a script and a recorded output in the code archive~\cite{Zenodo}, in the directories
\texttt{code/scope\_sumrule}, \texttt{code/sumrule\_freefield}, \texttt{code/hall\_sumrule\_cft}, \texttt{code/scope\_thermal} (with subdirectories
\texttt{hall} and \texttt{allorders}) and \texttt{code/no\_eft}; each directory has a \texttt{README.md} that lists the
scripts, their run times and their outputs.

\subsection{Conventions}
Take a CFT$_3$ in a homogeneous, isotropic thermal state at temperature $T$, with energy density $\varepsilon=2p$ and entropy
density $s$, and the shear components $T_1=(T_{xx}-T_{yy})/2$, $T_2=T_{xy}$ at zero spatial momentum. Rotation invariance
gives $G^R=\begin{pmatrix}g&-h\\h&g\end{pmatrix}$, with $G^R_{AB}(t)=i\theta(t)\langle[T_A(t),T_B(0)]\rangle$ integrated over
space. The circular combinations are $G_\pm=g\pm ih$, $\rho_\pm={\rm Im}\,G_\pm$, and $\delta$ denotes the thermal value minus
the vacuum value. At low frequency $g=p-i\eta\omega+O(\omega^2)$ and $h=-i\eta_H\omega+O(\omega^2)$. We define
\begin{equation}
\mathcal S_E=\frac1{\pi\varepsilon}\int_{-\infty}^{\infty}\frac{\delta\rho(\omega)}{\omega}\,d\omega,\qquad
\mathcal S_H=\frac1{\pi\varepsilon}\int_{-\infty}^{\infty}\eta_H(\omega)\,d\omega,\qquad
\eta_H(\omega)=-\frac{{\rm Im}\,\delta h(\omega)}{\omega},
\label{eq:tr-SESH}
\end{equation}
with $\rho={\rm Im}\,g$. On the gravity side we use the action
\begin{equation}
16\pi G\,\mathcal L=R+6-\tfrac16(\alpha_3J+\tilde\alpha_3Q)+\tfrac14(\alpha_4I^2+\alpha_4'P^2+2\tilde\alpha_4IP)+\dots,
\label{eq:tr-action}
\end{equation}
in units of the AdS radius, with $I=C_{\mu\nu\rho\sigma}C^{\mu\nu\rho\sigma}$ and $P=C_{\mu\nu\rho\sigma}\tilde C^{\mu\nu\rho\sigma}$,
and with $J=C_{\mu\nu}{}^{\rho\sigma}C_{\rho\sigma}{}^{\alpha\beta}C_{\alpha\beta}{}^{\mu\nu}$ and
$Q=C_{\mu\nu}{}^{\rho\sigma}C_{\rho\sigma}{}^{\alpha\beta}\tilde C_{\alpha\beta}{}^{\mu\nu}$. This is the action of
Ref.~\cite{CHLPS} with the Riemann tensor replaced by the Weyl tensor, which is exact on Einstein spaces up to Ricci and
topological terms. We write $\hat g_3=\alpha_3+i\tilde\alpha_3$, $G_4=\alpha_4+\alpha_4'=g_4/8\pi G$, and
$\mathcal V=(4\pi\eta/s-1)-i\,4\pi\eta_H/s$. The orientation is $\epsilon_{vrxy}=+r^2$ in the coordinates of
Sec.~\ref{sec:transport}.F; reversing it flips the sign of every Hall quantity together. The dictionary $t_4=-240\alpha_3$,
$|\alpha_T|=240|\tilde\alpha_3|$, i.e.\ $\tau=240|\hat g_3|$, is that of Sec.~\ref{sec:defs} [D].

\subsection{The shear sum rule}
\textbf{Statement} [Q; DV]. Under the conditions below, $\mathcal S_E=3/8+t_4/32$. This is the $d=3$ case of Eq.~(1.5) of
Ref.~\cite{Chowdhury:2016hjy} (restated with $t_2,t_4$ as their Eq.~(4.3)): the $t_2$ term drops out at $d=3$, and $\varepsilon=2p$.
Their $t_4$ is that of Refs.~\cite{Buchel,HofmanMaldacena}, the one used here (free scalar $t_4=4$, free fermion $t_4=-4$). The
left side is free of contact terms, since it equals $[\delta G(0)-\delta G(i\infty)]/\varepsilon$; the right side is $p$ from
hydrodynamics minus the $T$-exchange term of the OPE.

\textbf{Conditions} [Q\,+\,D].
(a)~$G^R$ analytic in the upper half plane, with $G^R(0)=p$ from hydrodynamics. At finite $c_T$ long-time tails make
$\eta(\omega)$ logarithmic as $\omega\to0$; this is integrable, so the sum rule survives.
(b)~No operator other than $T$ with $\Delta\le3$ in $T\times T$ has a thermal one-point function. In $d=3$ odd-spin one-point
functions vanish and unitarity puts other spinning operators above $\Delta=3$, so the condition is: no scalar or pseudoscalar
$\mathcal O$ in $T\times T$ with $\Delta_{\mathcal O}\le3$ and $\langle\mathcal O\rangle_\beta\ne0$. For $1<\Delta_{\mathcal O}<3$
the unsubtracted integral diverges; subtracting the tail $\propto\langle\mathcal O\rangle_\beta(i\omega)^{3-\Delta_{\mathcal O}}$
restores the sum rule with the same right side~\cite{Chowdhury:2016hjy,Lucas:2017dqa}. At $\Delta_{\mathcal O}=1$ the tail
coefficient is proportional to $d-\Delta_{\mathcal O}-2$ and vanishes~\cite{Lucas:2017dqa}. At $\Delta_{\mathcal O}=3$ the right
side can shift.
(c)~In a holographic CFT with large $c_T$, a single-trace scalar that couples to $TT$ (bulk $\phi W^2$) is sourced on the black
brane at the same order, so (b) amounts to: no single-trace scalar with $\Delta\le3$ couples to $TT$. This is implied by the
``no light single trace'' condition of Eq.~(5) of the Letter and is stronger than the condition $\Delta>1$ used for $\tau_*$.
The sum rule depends only on $t_4$ and $\varepsilon$: higher-derivative couplings that do not change $\langle TTT\rangle$,
including $R^4$, do not enter it, at any $c_T$.

\textbf{Checks} [DV].
\begin{itemize}
\item\emph{Free fields} (two independent codes: \texttt{scope\_sumrule/free\_fields\_sumrule.py} from Hamiltonian matrix
elements, \texttt{sumrule\_freefield/euclid\_tau.py} from the Euclidean path integral with symbolic Wick contractions and
explicit gamma traces). The vacuum correlators reproduce $c_T$ of Ref.~\cite{OsbornPetkou}, and $\varepsilon$ is taken from the
free energy ($\zeta(3)/\pi$ per real scalar, $3\zeta(3)/2\pi$ per Dirac fermion). The spectral function at zero momentum is a
pair continuum, $\rho_0=\omega^3/512$ times $\coth(\omega/4T)$ for the scalar and $\omega^3/256$ times $\tanh(\omega/4T)$ for the
Dirac fermion, plus a zero-frequency (Drude) weight, the free limit of the hydrodynamic peak, equal to $3\varepsilon/8$ in both
theories. The pair parts are $+\varepsilon/8$ and $-\varepsilon/8$. Hence $\mathcal S_E=1/2$ (scalar, $t_4=4$) and $1/4$
(fermion, $t_4=-4$), agreeing with $3/8+t_4/32$ to 20 digits. Without the Drude weight one would find $\pm1/8$. The improvement
term of the scalar stress tensor drops out at zero momentum; $\phi^2$ ($\Delta=1$) enters only through descendants that vanish at
zero momentum, and $\bar\psi\psi$ has no thermal one-point function in the free theory. In both theories the cross correlator
$\langle T_1T_2\rangle$ vanishes, so $\mathcal S_H=0$.
\item\emph{Einstein gravity} (\texttt{scope\_sumrule/einstein\_ads4\_sumrule.py}). On the AdS$_4$ black brane the integral gives
$\mathcal S_E=0.3749999601$; the integrand falls as $\omega^{-3.99}$, as the $\Delta=6$ double trace $[TT]$ predicts.
\item\emph{Large-$N$ Chern--Simons--fermion} (\texttt{sumrule\_freefield/cs\_fermion\_probe.py}), from the exact zero-momentum
thermal correlator of Ref.~\cite{Ghosh:2019sqf} and the energy density of Ref.~\cite{Aharony:2012ns}. The high-frequency constant
of the shear channel reproduces the $t_4/32$ dependence with $t_4=-4\cos\pi\lambda$ exactly (symbolic, at every $\lambda$; \texttt{hall\_sumrule\_cft/cs\_fermion\_exact.py}), up to a $\lambda$-independent
contact convention shared with the free Dirac fermion. This theory violates condition (b) at $\lambda\ne0$, through
$\bar\psi\psi$, so only the OPE side of the sum rule can be tested there.
\end{itemize}
Since the parity-even $\langle TTT\rangle$ of a CFT$_3$ is spanned by the free-scalar and free-fermion structures, $\mathcal S_E$
is affine in $t_4$, and the two free points fix it; the Einstein point is a third, independent check. We found no earlier independent check of the coefficient of $t_4$ (Ref.~\cite{Chowdhury:2016hjy} and its citing literature;
a search-level statement).

\subsection{The Hall sum rule}
\textbf{Theorem (Hall-channel sum rule)} [D]. Assume
(H1)~a unitary CFT$_3$ with a unique conserved spin-two operator $T$;
(H2)~the KMS state on $\mathbb R^2$ at temperature $T>0$ and zero chemical potentials, parity and time reversal possibly broken;
(H3)~the thermal OPE governs the large-Euclidean-frequency expansion of the plane-integrated $\langle T_AT_B\rangle_\beta$, the region
$|x|\gtrsim\beta$ contributing only terms that decay faster than any power;
(H4)~no scalar $\mathcal O$ in $T\times T$ with $\Delta_{\mathcal O}\le3$ has both $\langle\mathcal O\rangle_\beta\ne0$ and a nonzero parity-odd
$\langle TT\mathcal O\rangle$ coefficient (for $\Delta_{\mathcal O}<3$ it suffices that such terms be subtracted). Then
(1)~$\delta h(i\Omega)\to J_H$ as $\Omega\to+\infty$, with
\begin{equation}
J_H=\sigma\,\frac{\alpha_T\,\varepsilon}{32},\qquad \mathcal S_H:=\frac{J_H}\varepsilon=\sigma\frac{\alpha_T}{32},\qquad
\mathcal S_E-\tfrac38+i\sigma\mathcal S_H=\frac{t_4+i\alpha_T}{32},
\label{eq:tr-SHthm}
\end{equation}
where $\alpha_T$ is normalized by the collider energy matrix of Ref.~\cite{CDPrakash} and $\sigma=\pm1$ is the same for every CFT
(it depends only on orientation and helicity conventions); hence $\tau=32\,|\mathcal S_E-\tfrac38+i\sigma\mathcal S_H|$ exactly.
(2)~$J_H$ has no contact-term ambiguity. (3)~$\delta h(i\Omega)-J_H=O(\Omega^{3-\Delta_*})$, with $\Delta_*>3$ the lowest remaining
contributing dimension. (4)~If in addition (H5)~$\delta h^R-J_H$ is analytic in the upper half-plane and $O(|\omega|^{3-\Delta_*})$
uniformly up to the real axis, then $\mathcal S_H=(\pi\varepsilon)^{-1}\int_{\mathbb R}\eta_H(\omega)\,d\omega$, Eq.~\eqref{eq:tr-SESH},
with an absolutely convergent integral [D given H5].

\emph{Proof} (\texttt{code/hall\_sumrule\_cft/README.md}, Secs.~2--6). (i)~At zero chemical potential the rotation by $\pi$ in the
$(\tau,x)$ plane kills odd-spin one-point functions, and unitarity leaves only $T$ and scalars with $\Delta\le3$ as non-decaying terms
in the thermal OPE. (ii)~The most general $T$ term at zero spatial momentum is $J_E\delta_{AB}-{\rm sgn}(\Omega)J_H\epsilon_{AB}$. In
position space its even part is a $\delta(\tau)$ distribution, where contact terms live, so $J_E$ is scheme dependent and only
$\mathcal S_E$ is physical; its odd part is a separated-point $1/\tau$ tail, so $J_H$ is not. A rank count of $SO(3)$ invariants shows
that no parity-odd local term proportional to $\langle T\rangle_\beta$ exists; temperature-dependent parity-odd local terms arise only
from $\Delta=2$ scalars, and they are linear in $\omega$, not constant. (iii)~The OPE coefficient is the soft $|q|^3$ functional on the
soft leg of $\langle TTT\rangle$. (iv)~The homogeneous and parity-odd structures are supported only on equal helicities, and the
parity-odd one is the homogeneous one rotated by $\mathcal R_k$, $\mathcal R_k^2=-1$, on any one leg~\cite{Jain:2021gwa}; applied to a
hard leg it commutes with (iii). (v)~In the helicity basis the collider energy matrix of Ref.~\cite{CDPrakash} has flip element
$\propto t_4+i\alpha_T$, so the same real constant multiplies $t_4$ and $\alpha_T$. (vi)~That constant is fixed by the physical
$t_4$-slope of $\mathcal S_E$, which is $1/32$ (Ref.~\cite{Chowdhury:2016hjy}, with the checks of Sec.~\ref{sec:transport}.B).
(vii)~$\delta h^R(0)=0$ (the static response matrix is symmetric; equivalently $h_E$ is odd in $\omega_n$), and Kramers--Kronig under
(H5) gives the convergent form. Steps (iv) and (v) rest on the spinor-helicity results of Ref.~\cite{Jain:2021gwa}, read but not
re-derived. $\eta_H(\omega)$ is not sign-definite (in the bulk example below it is positive at $\omega=0$ while its integral is negative),
so the sum rule gives no bound by itself.

\emph{Checks} (both exact). (a)~Large-$N$ Chern--Simons--fermion theory: from the closed-form thermal correlator of
Ref.~\cite{Ghosh:2019sqf} and the gap equation, the large-frequency expansion of the helicity correlator is computed symbolically at
every 't~Hooft coupling $\lambda$; the Hall constant is exactly $\sin(\pi\lambda)/8=\alpha_T/32$ and the shear constant has slope
exactly $-t_4/32$, for $t_4=-4\cos\pi\lambda$, $\alpha_T=4\sin\pi\lambda$ (\texttt{hall\_sumrule\_cft/cs\_fermion\_exact.py}). The term
linear in frequency in that correlator is the exchange of $\bar\psi\psi$ ($\Delta=2$), phase-locked to the corresponding shear term,
not an arbitrary contact term; it must be subtracted before the integral form can be tested there. (b)~AdS$_4$ brane with the coupling
$\lambda_QQ$: the large-$\Omega$ constant is solved in closed form, $\mathcal S_H=-45\lambda_Q$ exactly, and the real-frequency integral
gives $-45.0002$ (\texttt{hall\_sumrule\_cft/bulk\_hall\_exact.py}); with $\tilde\alpha_3=-6\lambda_Q$ and $|\alpha_T|=240|\tilde\alpha_3|$
this is $|\alpha_T|/32$. The same calculation gives $4\pi\eta_H/s=108\lambda_Q$ as $\Omega\to0$ (Sec.~\ref{sec:transport}.F).

\textbf{Relative sign} [DV]. For the cubic parity-odd coupling, with both quantities taken from the same $h$,
$\mathcal S_H=-\tfrac5{12}\,(4\pi\eta_H/s)$ (numerically $-0.41659$ against $-5/12$): $h(i\Omega)$ starts at $+\eta_H\Omega$ and
ends at $J_H$ of the opposite sign. This ratio is independent of the orientation, since reversing it flips both. Equivalently,
in the helicity-$+$ channel $4\pi(\eta-i\eta_H)/s-1=18\hat g_3$ and $(\mathcal S_E-i\mathcal S_H)-3/8=-7.5\,\hat g_3$, so the two
are proportional with a real ratio. A structural argument gives the same result: both limits of $G_+$ are continuations of the
ingoing solution and depend on the single chiral coupling $\hat g_3$ [D-sketch].

\subsection{The universal disc}
\textbf{Statement} [D given Q; integral form of $\mathcal S_H$ D given H5]. Every unitary CFT$_3$ satisfying condition (b) of Sec.~\ref{sec:transport}.B and (H1)--(H4) of
Sec.~\ref{sec:transport}.C has $(\mathcal S_E-3/8)^2+\mathcal S_H^2\le1/64$. \emph{Proof.} The conformal-collider (energy positivity)
bound is $t_4^2+\alpha_T^2\le16$~\cite{CDPrakash}; by Eq.~\eqref{eq:tr-SHthm}, $(\mathcal S_E-3/8)^2+\mathcal S_H^2=\tau^2/1024$. The sign convention of $\mathcal S_H$ does not matter.
The parity-even projection, $1/4\le\mathcal S_E\le1/2$, is in Ref.~\cite{Chowdhury:2016hjy}. Free scalar and free fermion lie at
$(1/2,0)$ and $(1/4,0)$, mixtures of free fields on the diameter, Einstein gravity at the center, and the large-$N$
Chern--Simons--matter theories, with $t_4=-4\cos\pi\lambda$, $\alpha_T=4\sin\pi\lambda$~\cite{CDPrakash}, on the circle
for the high-frequency constants only: they violate condition (b) through $\bar\psi\psi$, so their integrals $\mathcal S_E$,
$\mathcal S_H$ need the subtraction of Sec.~\ref{sec:transport}.B before they can be placed on the disc. We found no earlier statement of the Hall sum rule or of the disc; the
searches covered the citing literature of Refs.~\cite{Chowdhury:2016hjy,CDPrakash} and quantum-Hall sum rules, so this is ``not
found'', not ``proven absent''.

\subsection{The spectral-to-gap theorem}
\textbf{Theorem SG} [D given the inputs listed]. Let $\delta_E=|\mathcal S_E-3/8|$ and
$\delta_{\rm th}=[(\mathcal S_E-3/8)^2+\mathcal S_H^2]^{1/2}$, both at $c_T=\infty$. Assume A1$'$, A2$'$, the inputs of
$\tau_*$ (Sec.~\ref{sec:angular}) and condition (b). Then
\begin{equation}
(\Delta_g-2)\Big[1-\frac{49}{(2\Delta_g+9)^2}\Big]^{1/2}\le\frac5{9\,\delta},\qquad
\Delta_g\le2.2401+\frac5{9\,\delta},
\label{eq:tr-SG}
\end{equation}
with $\delta=\delta_E$; with $\delta=\delta_{\rm th}$ it also uses the Hall sum rule, Eq.~\eqref{eq:tr-SHthm}, and hypothesis (H4).
\emph{Proof.} $\delta_E=|t_4|/32\le\tau/32$ and $\delta_{\rm th}=\tau/32$. Since $\tau_*$ is strictly decreasing, $\tau\le\tau_*(\Delta_g)$
is the first inequality. For the second write $x=\Delta_g-2$: the correction $x-x[1-49/(2x+13)^2]^{1/2}$ is at most $0.2401$,
attained at $\Delta_g\approx8.25$ (\texttt{scope\_sumrule/gap\_inverse.py}). $\square$\\
The theorem uses no black brane, no bulk effective theory and no positivity assumption. It is a statement at leading order in
$1/c_T$: the sum rule is exact at any $c_T$, but $\tau_*$ is not. Examples: $\delta=10^{-2}$ gives $\Delta_g\le57.64$ (exact
inversion) and $\delta=10^{-3}$ gives $\Delta_g\le557.8$.

\textbf{Fixed gap} [net grade of item (1): positivity certified, light remainder D-sketch; leading order in $1/c_T$]. If $\Delta_g\ge1000$ and no light single trace appears in
$T\times T$, $\delta_{\rm th}\le2.8914\times10^{-8}/32=9.04\times10^{-10}$; equivalently, $\delta_{\rm th}>9.04\times10^{-10}$
forces a higher-spin single trace below twist 1000.

\textbf{All-gap sharpening} [D given U1 and the inputs of Eq.~\eqref{eq:fp-allgap}]. For $\Delta_g\ge1000$ the analytic bound
$\tau\le280[\log\frac{\Delta_g+2}3]^{1/2}(\Delta_g-2)^{-2}$ gives
$\delta_{\rm th}\le f(\Delta_g):=\tfrac{35}4[\log\frac{\Delta_g+2}3]^{1/2}(\Delta_g-2)^{-2}$, decreasing in $\Delta_g$; so
$\delta_{\rm th}>f(D)$ with $D\ge1000$ forces a higher-spin state below twist $D$. Since $f(1000)=2.1\times10^{-5}$ exceeds the
fixed-gap value, this improves the inversion only for $\delta_{\rm th}<9.04\times10^{-10}$:

\begin{center}
\begin{tabular}{lcc}
$\delta_{\rm th}$ & bound from $\tau_*$ & bound from $f$\\\hline
$10^{-9}$ & $5.6\times10^{8}$ & below 1000 (fixed gap)\\
$10^{-10}$ & $5.6\times10^{9}$ & $5.5\times10^{5}$\\
$10^{-12}$ & $5.6\times10^{11}$ & $5.8\times10^{6}$\\
$10^{-14}$ & $5.6\times10^{13}$ & $6.0\times10^{7}$
\end{tabular}
\end{center}
The scaling is $\Delta_g\lesssim(8.75\sqrt{\log\Delta_g}/\delta_{\rm th})^{1/2}$.

\subsection{Shear and Hall viscosity of the brane}
\textbf{Setup.} Ingoing Eddington--Finkelstein coordinates,
$ds^2=-f\,dv^2+2\,dv\,dr+r^2(e^{H})_{ij}dx^idx^j$, $H=\begin{pmatrix}\phi&\chi\\\chi&-\phi\end{pmatrix}$, $f=r^2-1/r$, horizon at
$r=1$, with $\delta T^{xy}=-p\chi-\eta\dot\chi+\eta_H\dot\phi$. On the background $I=12/r^6$, $J=12/r^9$, $P=Q=0$. All four
invariants have vanishing first variation under the homogeneous shear, because a spin-two perturbation has no rotationally
invariant linear image; at the horizon $I=J=12$, $I'=-72$, $J'=-108$, $I''=504$, $J''=1080$.

\textbf{Shear} [DV]. For each operator $\mathcal O$ in $16\pi G\mathcal L=R+6+\lambda\mathcal O$, the static shear mode and the Wald
entropy give $\delta(4\pi\eta/s)$ per unit $\lambda$ (\texttt{scope\_thermal/etas\_diag.py}, \texttt{thermo\_check.py}):
$I$: $0$; $J$: $-108$; $I^2$: $+864$; $IJ$: $-1944$; $J^2$: $-6480$; $I^3$: $+31104$; $P$, $Q$, $P^2$, $IP$: $0$. The $J$
coefficient is $-(3/40)t_4$ with $t_4=-240\alpha_3$, as in Ref.~\cite{Li:2019auk}; the $I^2$ coefficient agrees with the M-theory
result of Ref.~\cite{Garousi:2008ai} after converting conventions; the Wald entropy agrees with the on-shell free energy through
the first law for $J$ and $I^2$. Parity-odd operators cannot shift $\eta$ at linear order: the brane and the static shear are
invariant under $x\leftrightarrow y$.

\textbf{Hall} [DV]. Keeping both polarizations and one time derivative, the mixed quadratic terms are
\begin{equation}
P^{(2)}\supset-\frac6{r^2}\dot\phi\,\chi'',\qquad Q^{(2)}\supset\frac{54}{r^6}\dot\phi\,\chi'+\frac{18}{r^5}\dot\phi\,\chi''.
\end{equation}
Terms without a time derivative vanish (a static configuration is invariant under the orientation-reversing $t\to-t$), and no
term involves $\phi$ without a time derivative. The $\chi$ equation at $O(\omega)$ is $\partial_r\Pi=\lambda\sigma(r)\dot\phi$ with
$\Pi=-r^2(f\chi'+\dot\chi)$, regularity at the horizon fixes $\Pi(1)=-\dot\chi$, and the curvature terms do not reach the boundary
since $C\sim r^{-3}$. Hence $4\pi\eta_H/s=\lambda\int_1^\infty\sigma\,dr$, with
$\sigma_Q=432/r^5$ and $\sigma_{IP}=-3024/r^8$:
\begin{equation}
\delta(4\pi\eta_H/s)=+108\ (Q),\qquad -432\ (IP),\qquad 0\ (P,\ P^2,\ I,\ J,\ I^2).
\end{equation}
The odd terms do not change the background or the Wald entropy at linear order. These numbers were obtained twice, by a
symbolic calculation in the bivector representation of the Weyl tensor and by an explicit index computation
(\texttt{scope\_thermal/hall/hall\_ef.py}). Checks: (a)~the response of the $\phi$ channel to $\dot\chi$ is $-108$ and $+432$,
so the response matrix is antisymmetric, as Onsager reciprocity requires; (b)~a coupling $\theta(r)P$ gives $\sigma=-6\theta''$
and $6\theta'(1)$, the gravitational Chern--Simons result of Ref.~\cite{Saremi:2011ab} up to the conversion of conventions, and
$IP$ is exactly $\theta P$ with $\theta=12/r^6$, so this checks the $IP$ coefficient directly; (c)~a constant coefficient of
$P$ gives zero, as it must for a topological term, while $Q$ is not a total derivative. In terms of Eq.~\eqref{eq:tr-action},
\begin{equation}
\mathcal V=18\,\hat g_3+216\,(\alpha_4+i\tilde\alpha_4)+\dots,\qquad 4\pi\frac{|\eta_H|}s=\frac3{40}|\alpha_T|\ \text{(cubic order)}.
\label{eq:tr-V}
\end{equation}
Hall viscosity is a frame-invariant, dissipationless first-order transport coefficient of a parity-violating fluid in
$2+1$ dimensions~\cite{Jensen:2011xb}; at zero momentum the only local parity-odd contact term is the gravitational Chern--Simons
term, of order $\omega^3$, so the coefficient of $\omega$ is free of scheme ambiguity. Earlier holographic Hall viscosities
required an axion with a radial profile or a magnetic field~\cite{Saremi:2011ab,Hoyos:2019pyz}; we found no axion-free
computation for a neutral brane and no relation to $\alpha_T$ in the literature.

\textbf{Quartic positivity} [D; DV]. For Eq.~\eqref{eq:tr-action} the flat-space four-graviton contact amplitudes are: all-plus
$\propto\alpha_4-\alpha_4'-2i\tilde\alpha_4$ and maximally helicity violating $\propto\alpha_4+\alpha_4'$, with $IP$ not
contributing to the latter (script \texttt{quartic\_IP.py} in \texttt{scope\_thermal/hall}; the normalization ratio to $\tfrac12\hat g_4(s^2+t^2+u^2)^2$
is $1.000000$). The flat-space bound $|\hat g_4|\le g_4$~\cite{CHLPS} is then $(\alpha_4-\alpha_4')^2+4\tilde\alpha_4^2\le G_4^2$,
i.e.\ $\alpha_4,\alpha_4'\ge0$ and $\tilde\alpha_4^2\le\alpha_4\alpha_4'$. Our parity-odd normalization differs by a factor of 2
from the one implied in Ref.~\cite{CHLPS}; we use the computed amplitude. Consequently $|\alpha_4+i\tilde\alpha_4|^2\le\alpha_4G_4\le G_4^2$.
This is a flat-space input; we have not established it in AdS.

\textbf{Numbers at $\Delta_g=1000$} [D given the inputs]. $\Lambda=998.50$. Cubic: $(3/40)\times2.8914\times10^{-8}=2.1686\times10^{-9}$.
Quartic: with quartic positivity and $\tilde g_4=g_4$ (L7), $G_4\le92.81/\Lambda^6=9.36\times10^{-17}$ and $216G_4=2.0\times10^{-14}$
($3.6\times10^{-14}$ without the scalar condition). Second order: $(17/3840)\tau^2=3.7\times10^{-18}$. Remainder (Sec.~\ref{sec:transport}.H):
$1.95\times10^{-20}C$. Total $|\mathcal V|\le2.17\times10^{-9}$ and $|\eta_H|/s\le1.73\times10^{-10}$. Without positivity, the cubic
part alone is at most $(3/40)\tau_*(1000)=1.34\times10^{-3}$, i.e.\ $|\eta_H|/s\le1.06\times10^{-4}$.

\subsection{General linear response}
\textbf{Algebraic invariants} [DV]. For $16\pi G\mathcal L=R+6+F(I,J,P,Q)$, at first order in $F$, with derivatives at
$(I,J,P,Q)=(12,12,0,0)$,
\begin{align}
\delta(4\pi\eta/s)&=432F_{II}-648F_{IJ}-1944F_{JJ}-108F_J,\\
\delta(4\pi\eta_H/s)&=108F_Q+1296F_{QI}+1944F_{QJ}-432F_{PI}-648F_{PJ}.
\end{align}
These follow from the response to generic radial profiles $\theta_i=\partial F/\partial y_i$, $y=(I,J,P,Q)$, which we computed
directly (\texttt{scope\_thermal/allorders/shear\_theta.py}, \texttt{hall\_theta.py}): $\delta(4\pi\eta/s)=-6\theta_I'(1)-108\theta_J(1)+18\theta_J'(1)$
and $\delta(4\pi\eta_H/s)=6\theta_P'(1)+108\theta_Q(1)-18\theta_Q'(1)$; for $Q$ the source is an exact total derivative, so the
result is horizon-local. No $F_P$, $F_{PP}$, $F_{PQ}$ or $F_{QQ}$ term can appear. Monomial checks against direct computation:
$IQ$ $+2592$, $JP$ $-648$, $JQ$ $+3240$, $I^2P$ $-10368$, $IJ$ $-1944$, $I^3$ $+31104$, $J^2$ $-6480$.

\textbf{Gradient extension} [D; DV]. For $F(y_i,X_{ij})$ with $X_{ij}=\nabla y_i\cdot\nabla y_j$, define the Euler derivatives
$A_i=\partial F/\partial y_i-\nabla_\mu[\partial F/\partial(\nabla_\mu y_i)]$ on the brane. Then
$\delta(4\pi\eta/s)=-108A_J-6A_I'+18A_J'$ and $\delta(4\pi\eta_H/s)=108A_Q+6A_P'-18A_Q'$. The shear leaves the $(v,r)$
inverse metric and $\sqrt{-g}$ unchanged and the invariants have no first variation, so $X^{(2)}=2\nabla y^{(0)}\cdot\nabla y^{(2)}$
and integration by parts leaves the Euler derivatives; the extra Iyer--Wald term is proportional to $\nabla y$, which vanishes on
the bifurcation surface. Direct checks (\texttt{gradient\_direct.py}, \texttt{hall\_gradient\_direct.py}): $(\nabla I)^2$ gives
$31104$ (flux $32832$ minus Wald $1728$), $\nabla I\cdot\nabla P$ gives $-15552$ (time-dependent computation with no integration
by parts by hand), $(\nabla J)^2$ gives $-279936$.

\textbf{Ten derivatives} [DV\,+\,Q]. With $U=(\nabla C)^2$ and $\tilde U=\nabla C\cdot\nabla\tilde C$, on any Einstein space with
$R_{\mu\nu}=-3g_{\mu\nu}$, $U=\tfrac12\Box I+6I+3J$ and $\tilde U=\tfrac12\Box P+6P+3Q$; on the brane $U=180f/r^8$ and
$\Box I=216/r^6-432/r^9$. The seven independent ten-derivative structures~\cite{Ruhdorfer:2019qmk} are $IJ$, $IQ$, $PJ$, $PQ$,
$IU$, $PU$, $P\tilde U$, with
\begin{center}
\begin{tabular}{lccccccc}
& $IJ$ & $IQ$ & $PJ$ & $PQ$ & $IU$ & $PU$ & $P\tilde U$\\\hline
$\delta(4\pi\eta/s)$ & $-1944$ & $0$ & $0$ & $0$ & $-16200$ & $0$ & $0$\\
$\delta(4\pi\eta_H/s)$ & $0$ & $2592$ & $-648$ & $0$ & $0$ & $3240$ & $0$
\end{tabular}
\end{center}
$IU$ was also computed without the identity (flux $-15768$ plus Iyer--Wald $-432$; \texttt{boxC\_wald.py}), and $PU$ is $\theta P$
with $\theta=U$. The count is a flat-space Hilbert-series count; with a cosmological constant each reduction also produces
lower-dimension terms times powers of the curvature scale, which renormalize $t_4$ exactly and the quartic coupling measured by
$\langle TTTT\rangle$ up to $O(\Lambda^{-2})$.

\subsection{Higher-derivative remainder}
\textbf{Proposition R} [D; DV, for the stated class and basis]. Let $x=\Lambda^{-2}\le1/997002$ ($\Delta_g\ge1000$). Consider the
class consisting of all algebraic monomials $I^aJ^bP^cQ^d$, the ten-derivative structures $IU$, $PU$, $P\tilde U$, and all monomials
of $F(y,X)$ at twelve or more derivatives, and suppose every coefficient $\gamma$ of an operator with $2k$ derivatives, $k\ge5$,
obeys $|\gamma|\le C\,x^{k-1}$ in this basis. Then at first order in the couplings their total contribution obeys
$|R_E|\le18145\,C/\Lambda^8$ and $|R_H|\le6481\,C/\Lambda^8$, so $|R|\le19300\,C/\Lambda^8=1.95\times10^{-20}C$ at $\Delta_g=1000$.
\emph{Proof.} With $A=12x^2$, $B=12x^3$, the sums of the absolute term-wise responses of the algebraic tower are
$E(x)=\frac{864x^4}{(1-A)^3(1-B)}+\frac{648x^5}{(1-A)^2(1-B)^2}+\frac{3888x^6}{(1-A)(1-B)^3}+\frac{108x^3}{(1-A)(1-B)^2}$ and
$H(x)=\frac{432x^4}{(1-A)^2(1-B)}+\frac{648x^5}{(1-A)(1-B)^2}+\frac{108x^3}{(1-A)(1-B)}+\frac{1296x^5}{(1-A)^2(1-B)}+\frac{1944x^6}{(1-A)(1-B)^2}$
(they agree exactly with brute-force sums over monomials). After the cubic and quartic terms are removed, their normalized
remainders at the largest $x$ are $1944.0377$ and $3240.0136$ in exact rational arithmetic. Gradient monomials with three or
more factors of $X$ give zero at the horizon; those with one or two give, summed exactly, at most $34993x$ and $104977x$ beyond
ten derivatives. Adding $IU$ ($16200$) and $PU$ ($3240$) gives $18144.2$ and $6480.1$ (\texttt{allorders/remainder.py}). $\square$\\
\emph{Scope.} The constants depend on the basis: in the basis $\{(\nabla I)^2,\nabla I\cdot\nabla P\}$ instead of $\{IU,PU\}$ the
ten-derivative constants become $33048$ and $18792$. Beyond ten derivatives the class does not contain every operator, for example
$(\Box I)^2$ or general $C^4\nabla^4$ structures, and the hypothesis on $\gamma$ is an assumption about couplings of five or more
gravitons, which the vacuum three- and four-point functions do not constrain (Sec.~\ref{sec:transport}.K). Proposition R is
therefore not an all-orders statement about the bulk effective theory. Numerically, with natural coefficients
$C\lesssim100$, the remainder is below $2\times10^{-18}$.

\subsection{Second order in the cubic couplings}
\textbf{Result} [DV]. At tree level, $4\pi\eta/s=1+18\alpha_3-255\alpha_3^2=1-\tfrac3{40}t_4-\tfrac{17}{3840}t_4^2$, which reproduces
the $t_4^2$ coefficient of Ref.~\cite{Li:2019auk}, and $4\pi\eta_H/s=-18\tilde\alpha_3+255\alpha_3\tilde\alpha_3$; $\eta$ has no
$\tilde\alpha_3^2$ term and $\eta_H$ no $\alpha_3^2$ or $\tilde\alpha_3^2$ term. Hence $\mathcal V=\hat g_3(18-255\,{\rm Re}\,\hat g_3)+\dots$
and $|\mathcal V-18\hat g_3|\le255|\hat g_3|^2=\tfrac{17}{3840}\tau^2$ (\texttt{allorders/second\_order\_shear.py},
\texttt{hall\_second\_order.py}, \texttt{wald\_J\_corrected.py}). The Einstein kinetic coefficient is $-Nr^2f/2$ on any static
background, so the second-order background is not needed. Mixed cubic--quartic terms were not computed; under the bounds they
are below $10^{-24}$.

\subsection{The corrected viscosity bound and the transport disk}
Combine Eq.~\eqref{eq:tr-V}, the exact sum rules, and the relative sign of Sec.~\ref{sec:transport}.C:
\begin{equation}
\mathcal Z:=\mathcal V+\tfrac{12}5\big(\mathcal S_E-\tfrac38-i\mathcal S_H\big)=216(\alpha_4+i\tilde\alpha_4)-255\,({\rm Re}\,\hat g_3)\,\hat g_3+R,
\end{equation}
since $\tfrac{12}5\cdot\tfrac1{32}=\tfrac3{40}$ and the cubic parts cancel in both components. With the combination
$\mathcal S_E-3/8+i\mathcal S_H$ instead, the cubic imaginary parts add. Taking the real part and $\alpha_4\ge0$:
\begin{equation}
4\pi\frac\eta s\ge1-\frac{12}5\Big(\mathcal S_E-\frac38\Big)-\frac{17}{3840}t_4^2-|R_E|
\end{equation}
[D given the bulk effective theory on the brane, Proposition R's hypothesis, and $\alpha_4\ge0$]. With
$\tilde\alpha_4^2\le\alpha_4\alpha_4'$, $|\mathcal Z-108G_4|\le108G_4+\tfrac{17}{3840}\tau^2+|R|$: a disc tangent to the origin,
so that a corrected shear response at either end of its allowed interval forces the corrected Hall response to vanish. The
combination is not exact in the CFT: $\eta/s$ needs the brane. A single-trace scalar with $\phi W^2$ coupling and
$3<\Delta<\Delta_g$ shifts $\eta/s$ without shifting $\mathcal S_E$ and must be excluded, or controlled as follows.
\emph{Induced scalar hair} [D]. On the thermal branch continuously connected to AdS$_4$-Schwarzschild with $\phi=0$, a coupling
$\lambda\phi C^2$ sources $\phi$ at order $\lambda$. The static source-free solution is $\phi_{\rm sf}=Bz^a{}_2F_1(a,a;2a;z)$, with
$z=(r_h/r)^3$ and $a=\Delta/3$, and near the horizon ${}_2F_1(a,a;2a;z)=-[\Gamma(2a)/\Gamma(a)^2]\log(1-z)+O(1)$; since
$\Gamma(2a)/\Gamma(a)^2\ne0$ for $a>0$, regularity forces $B=0$, so with the usual boundary conditions the regular solution has
$\phi=O(\lambda)$. The scalar stress and the explicit metric correction then start at $O(\phi^2)$ and $O(\lambda\phi)$, so
$\delta(4\pi\eta/s)=O(\lambda^2)$, which is $O(\Delta_g^{-4})$ if $\lambda=O(\Delta_g^{-2})$. This replaces the exclusion by a
power-counting assumption of the same kind as Proposition R; it does not cover disconnected hairy phases.

\subsection{What an EFT-free derivation of \texorpdfstring{$\eta/s$}{eta/s} requires}
The transport statements above use the bulk effective theory on the brane. A CFT derivation of $\eta/s=1/4\pi$ up to powers
of $1/\Delta_g$ remains open. We record what it requires.

\textbf{Remark 1 (order of limits)} [Q]. The target is $\mathcal V_{\rm pl}=\lim_{\omega/T\to0^+}\lim_{c_T\to\infty}4\pi\eta(\omega)/s$.
In $2+1$ dimensions long-time tails give the leading-log running $\mu\,dg/d\mu=-1/(16\pi c_{\rm th}g)$, with $c_{\rm th}=s/T^2$ and
$g=\eta/s$~\cite{Kovtun:2012rj}. Resolving a deviation of order $\Delta_g^{-4}$ at finite $c_T$ therefore needs
$\log(\mu_0/\omega)\ll c_T/\Delta_g^4$.

\textbf{Remark 2 (multi-stress data)} [D-sketch]. Expand the thermal $\langle TT\rangle$ in the thermal OPE over $T$ and the
multi-stress tensors $[T^n]$. At leading order every $n$ contributes to the connected $O(c_T)$ correlator. $n=1$ is fixed by
$c_T$, $t_4$ and $\alpha_T$. At $n=2$ the contribution is $c_T[(\lambda_1b_0+\lambda_0b_1)B_6+\lambda_0b_0\gamma_1\partial_\Delta B_6]$,
schematically for the lowest family: the $O(1/c_T)$ OPE data $\lambda_1,\gamma_1$ come from $\langle TTTT\rangle$, but the
subleading thermal one-point coefficient $b_1$ of $[TT]$ is independent data (cf.\ Eq.~(3.17) of Ref.~\cite{Karlsson:2022osn}
in four dimensions). $n\ge3$ needs higher-point vacuum data $\langle TT[T^n]\rangle$ and further thermal one-point data. What
must be bounded is the deviation from the Einstein hierarchy, which is itself nonzero at infinite gap. Finally the thermal OPE
converges only at separations below $\beta$, while $\eta$ is a zero-frequency quantity, so control of the real-time continuation
is also needed.

\textbf{Remark 3 (type-N blindness)} [D]. If the bivector Weyl operator $W_0$ (the self-dual $3\times3$ block) satisfies
$W_0^2=0$, then $F=I_2I_3={\rm tr}\,W^2\,{\rm tr}\,W^3$ obeys $F=DF=D^2F=0$ at $W_0$. \emph{Proof.} $I_2=I_3=0$ at $W_0$, and
$DI_3[X]=3\,{\rm tr}(W_0^2X)=0$; the product rule gives $DF=DI_2\,I_3+I_2\,DI_3=0$ and
$D^2F=D^2I_2\,I_3+2\,DI_2\,DI_3+I_2\,D^2I_3=0$. $\square$ The same holds for $\sqrt{-g}F(W(g))$. Aligned AdS plane and shock waves
of any profile, including superpositions of parallel waves, are of type N, so their quadratic fluctuations are blind to $I_2I_3$.
This bears on causality constraints from shock waves~\cite{Belin:2019mnx} and on detector methods. It is sharp at cubic order:
for $N=\begin{pmatrix}1&i&0\\i&-1&0\\0&0&0\end{pmatrix}$ ($N^2=0$) and $X={\rm diag}(2,1,-3)$,
${\rm tr}(N+tX)^2\,{\rm tr}(N+tX)^3=18t^3+90t^4-252t^5$.

\textbf{Remark 4 (no closure from four-point data)} [D]. Integrating out a scalar of mass $M$ coupled to $g_2I_2+g_3I_3$ gives
$aI_2^2+bI_2I_3+cI_3^2$ with $a=g_2^2/2M^2$, $b=g_2g_3/M^2$, $c=g_3^2/2M^2$, and sums of such terms obey $b^2\le4ac$. For any
$A>0$ and any $B$, $g_2=\sqrt{2A}M$, $g_3=BM/\sqrt{2A}$ give $a=A$, $b=B$, $c=B^2/4A$. A bound on the quartic coefficient plus
mixed positivity therefore does not bound the quintic coefficient $b$; an independent bound on the six-point diagonal $c$ is
needed.

\textbf{Remark 5 (spectral deformations)} [D]. Sum rules, positivity and the KMS condition do not fix $\eta$. Equal
deformations $\delta\rho_\pm=\lambda\omega(e^{-\omega/u}-\tfrac uve^{-\omega/v})$ preserve $\mathcal S_E$ and the OPE but change
$\eta$ by $\lambda(1-u/v)$; opposite ones, $\delta\rho_\pm=\pm\lambda\omega^2e^{-\omega/u}$, preserve $\mathcal S_H$ and change
the infrared Hall moment by $2\lambda u/\pi$. More sharply, $\delta\rho_\epsilon=A\omega e^{-\omega^2/\epsilon^2}P_m(\omega^2/\epsilon^2)$
with $P_m=L_m^{(-1/2)}(x)/L_m^{(-1/2)}(0)$ has $\delta\eta=A$, vanishing moments $M_0,\dots,M_{m-1}$, and
$\sup_\tau|\delta G_E|\le2A\,m!\,\zeta(2m)\beta^{2m-1}\epsilon^{2m+1}/[\sqrt\pi(2\pi)^{2m}]$. For $m=1$,
$\delta\rho=A\omega(1-2\omega^2/\epsilon^2)e^{-\omega^2/\epsilon^2}$, $\sup|\delta G_E|\le A\beta\epsilon^3/(12\sqrt\pi)$, and positivity
is preserved if $\rho_0\ge\kappa\omega$ and $A\le\kappa e^{3/2}/2$. Its time kernel has $\int\delta K=0$ and
$\int|\delta K|=4A\epsilon/(e\sqrt\pi)\to0$, while $-\int t\,\delta K=A$ stays fixed. This is a response-function construction,
not a CFT: thermal crossing or higher-point consistency would have to exclude it.

\textbf{Remark 6 (a sufficient Kubo lemma)} [D]. Write the subtracted correlator as $g-g_0=h_{\rm UV}+\int(e^{i\omega u}-1)\,\delta k(u)\,du$,
where $g_0$ is the Einstein reference with $\eta_0=s_0/4\pi$, and suppose $\int|u\,\delta k|\le E$,
$|\lim_{\omega\to0}{\rm Im}\,h_{\rm UV}/\omega|\le U$ and $|s-s_0|\le\delta s_0$. Then $|4\pi\eta/s-1|\le[4\pi(E+U)/s_0+\delta]/(1-\delta)$.
A uniform $\Delta_g^{-4}$ bound on the resummed first time moment would therefore complete the Kubo step.

\textbf{Remark 7 (the viscosity shift is a cubic amplitude on the horizon curvature)} [D]. For $S=\frac1{16\pi G}\int\sqrt{-g}\,[R+6+\lambda h(I,J)]$
with $h$ any function of the full-tensor invariants $I=C^2$, $J=C^3$, the general response formula above gives
$4\pi\eta/s=1-108\lambda\Gamma_h+O(\lambda^2)$, $\Gamma_h=[h_J-4h_{II}+6h_{IJ}+18h_{JJ}]_{I=J=12}$. Write the chiral Weyl blocks on the type-D curve
$W_\pm=q\,{\rm diag}(-2,1,1)$, so $I=48q^2$, $J=-96q^3$ and the horizon is $q=-\tfrac12$. For the null polarizations $P(z)=nn^T$,
$n=(2z,1-z^2,i(1+z^2))$, $U=P(0)$, $V=\bar U$ ($P^2=U^2=V^2=0$), which are the chiral Weyl tensors of three conserved on-shell momenta
$\lambda\tilde\lambda_h$, $-\lambda(\tilde\lambda_h+z\tilde\lambda_s)$, $z\lambda\tilde\lambda_s$, the third variation of $h$ in the $+$ block is
$D^3h[P(z),U,V]=-384z^2\,\Gamma(q)$ with $\Gamma(q)=h_J+8qh_{II}+24q^2h_{IJ}-144q^3h_{JJ}$ and $\Gamma(-\tfrac12)=\Gamma_h$ [verified symbolically for
$h=J,IJ,I^2,J^2$]. So the viscosity shift is $-108$ times the coefficient, relative to $C^3$, of a six-derivative cubic principal amplitude on the
horizon curvature. For $h=\beta_5IJ$, $\Gamma(q)=72\beta_5q^2$: it vanishes on the vacuum ($q=0$), which is Remark~3 in amplitude form, and equals
$18\beta_5$ at the horizon. A CFT argument would have to bound this background-dependent amplitude; the match to a full eikonal time delay
(exchanged leg, propagation, derivative operators) is not done. This single coefficient is not the whole target: the five-graviton gradient operator $Y=\nabla I\cdot\nabla J$, with
nonzero flat principal amplitude $\propto s_{12}\langle12\rangle^4[34]^2[45]^2[53]^2$, has vanishing projection on this frozen-background
channel (conserved three-point kinematics kill the gradient terms), yet shifts $4\pi\eta/s$ by $-34992\lambda$, from the gradient-extended formula
$4\pi\eta/s-1=\lambda[-6\mathcal A_I'(1)+18\mathcal A_J'(1)-108\mathcal A_J(1)]$ with Euler derivatives $\mathcal A_{I,J}$ [D; $(\mathcal A_I,\mathcal A_J)_h=(324,216)$,
$(\mathcal A_I',\mathcal A_J')_h=(-5832,-2592)$]. Background gradients must therefore also be controlled.

\textbf{Remark 7$'$ (positive subtraction removes the interference)} [D]. If a positive functional $L(M)={\rm tr}(WM)$, $W\succeq0$, annihilates a
spectator Gram matrix $A=BB^\dagger$, then $WB=0$ and $L(BC^\dagger+CB^\dagger)=0$ for every new amplitude $C$. For $X=[T_1,[T_2,T_3]]=gX_1+g^2X_2+g^3X_3$,
$g=c_T^{-1/2}$, removing $\langle X_1^\dagger QX_1\rangle$ therefore also removes the $X_1$--$X_3$ interference through which a quintic vertex enters linearly.
Positivity of higher moments plus a gap does not bound the higher diagonal either ($O=t\mathbf 1$ keeps $R_4$ fixed while $R_6=t^2R_4$ grows).
These obstruct this construction on the full matrix cone, not every bootstrap functional. Nor does a bounded one-particle ANEC matrix element control the whole quintic coefficient: a heavy scalar $X$ coupled to
$g_2I+g_3J$ gives $\beta_5=g_2g_3/m^2$ through a \emph{single} exchange (stripped amplitude $g_2g_3/(m^2-s_{12})+c_5$, with no double residue), and
$g_3$ is $XTTT$ data absent from $X_aTX_b$. In the five-point comb these single-trace--double-trace and double-trace--double-trace sectors enter at the
same planar order as the double single-trace sector, and the Cauchy decomposition of a two-variable dispersion keeps them in its single-variable
subtraction functions. A Gram bound $|B_5|^2\le M_4M_6$ needs a six-point diagonal, which the scalar example saturates for every $g_3$
($b_5^2=4a_4a_6$).

\textbf{Remark 8 (uniform Kubo limit)} [D]. Let the zero-momentum shear action be $-\tfrac{\mathcal N}2\int[A\phi'^2-B\dot\phi^2]$ with
$A=f_0a/u^2$, $B=b/(u^2f_0)$, $f_0=1-u^3$, and $\max(\|a-1\|_\infty,\|b-1\|_\infty,\|c'\|_\infty)\le\varepsilon\le\tfrac14$, $c=\sqrt{ab}$. Stripping
the ingoing factor turns the radial equation into a Volterra equation whose kernel has norm at most $|\omega|M_*$, with
$M_*=\sqrt{5/3}\,(\tfrac13+\tfrac23M_0)=1.1633$, $M_0=\tfrac12\log3+\pi/(6\sqrt3)$ (the Einstein value). Hence, for $|\omega|M_*\le\tfrac14$,
$|{-{\rm Im}\,G_R}/(\omega\mathcal Nc(1))-1|\le7M_*|\omega|$, uniformly in the class, and $\eta=\mathcal Nc(1)$ with $|\eta/\mathcal N-1|\le\varepsilon$.
With $s=4\pi\mathcal N(1+\nu)$, $|\nu|\le\varepsilon$: $|4\pi\eta/s-1|\le2\varepsilon/(1-\varepsilon)$. The horizon logarithm is kept exactly, so the
$\omega\to0$ limit is uniform; what remains is to derive the action class and $\varepsilon=O(\Delta_g^{-4})$ from the CFT. Note also that the
collider estimate of Ref.~\cite{AKT} carries $\log\Delta_g/\Delta_g^4$, so a strict $\Delta_g^{-4}$ rate would need a sharper argument.

\textbf{Remark 9 (a relaxation gap does not control the Kubo limit)} [D]. Let $z=-i\omega$, $F_M(z)=\prod_{j=1}^M\gamma_j/(z+\gamma_j)$ with
$\gamma_j\in[\gamma,\Gamma)$, $G_0=B\Gamma z/(z+\Gamma)$ and $G_M=G_0+AzF_M$, $B>A>0$. Every pole stays at distance $\ge\gamma$ below the real axis, the
spectral density obeys $\rho_M\ge2(B-A)\Gamma^2\omega/(\Gamma^2+\omega^2)>0$ (so a KMS two-point completion exists), the first $M-2$ inverse-frequency
coefficients agree with $G_0$, and $\sup_\tau|\delta G_E(\tau)|\le A\Gamma^3q^{M-3}/(12T)$ with $q=\Gamma/(\Gamma+2\pi T)<1$; yet the Kubo slopes differ
by exactly $A$. The causal kernel is $\delta K_M=Ap_M'$ with $p_M$ the density of a sum of $M$ exponential times, and
$\int t|\delta K_M|\,dt\gtrsim A\sqrt M$. So a pole gap plus exponentially close Euclidean data still leaves a uniform analytic norm on a disk
$|\omega|\le r$ (or a weighted causal bound) to be proved; with such a bound of size $K\Delta_g^{-4}s/4\pi$ and an entropy bound, Cauchy's estimate
gives the theorem. These are response functions, not CFT correlators.

\textbf{Remark 10 (relative form bounds suffice, without a gap)} [P, conditional]. With the Kubo--Mori product, $Q=[H,\cdot]^2\ge0$ and the
entropy-normalised regulated shear source $a$, $\eta/s=\lim_{\nu\downarrow0}\langle a,\nu(Q+\nu^2)^{-1}a\rangle$. If a reference and the theory share a
Hilbert space with $(1-\epsilon)Q_0\le Q_1\le(1+\epsilon)Q_0$ and $\langle\delta a,R_0\delta a\rangle\le\delta^2\langle a_0,R_0a_0\rangle$, $R_0=\nu(Q_0+\nu^2)^{-1}$,
operator inverse order and $\nu(cQ_0+\nu^2)^{-1}=c^{-1/2}R_0(\nu/\sqrt c)$ give $|4\pi\eta/s-1|\le2\epsilon+4\delta$ ($\epsilon,\delta\le\tfrac12$, reference value
$1/4\pi$). No spectral gap is needed. The open step is to derive $\epsilon,\delta=O(\Delta_g^{-4})$: ${\rm spec}\,[H,\cdot]$ consists of energy
differences, which the vacuum gap does not control.

\textbf{Remark 11 (Euclidean data give at most a logarithmic rate)} [P]. Let $E(s)=4\pi[\delta G_R(0)-\delta G_R(2\pi iTs)]/(2\pi Ts\,s_{\rm ent})$, the
Poisson average of $4\pi\delta\eta/s$ at width $2\pi Ts$, so $E(0^+)=4\pi\eta/s$ and $E(n)$, $n\ge1$, is a Euclidean $S^1\times\mathbb R^2$ quantity free of contact
terms [Einstein: $E(1..4)=0.41275,0.23820,0.16375,0.12400$, checked to $10^{-10}$]. If $F=E-E_{\rm Einstein}$ is analytic with $|F|\le M$ on ${\rm Re}\,s>-\gamma$
and $|F(n)|\le\varepsilon$ for $n\le N$, then $|4\pi\eta/s-1|\le C_\gamma M(\log M/\varepsilon)^{-2\gamma}$ with explicit constants (half-plane Blaschke
products and Lagrange interpolation); an explicit family that keeps $\rho>0$ and every OPE coefficient and sum rule shows this rate is sharp to a factor
2--3. In the bulk the reason is that $\eta$ is horizon-local while the $n$-th Matsubara mode vanishes as $\rho^{|n|}$ at the tip of the cigar: a change
within depth $w$ of the horizon shifts $4\pi\eta/s$ by a fixed amount while $F(n)\propto w^n$. Euclidean thermal data therefore cannot give a power of
$\Delta_g$; a power needs a small real-time norm (Remark 10) or structural input, e.g.\ a sign-definite relaxation-function deviation supported at
$t\le t_L$, which gives $|F(0)|\le e^{2\pi Tt_L}|F(1)|$ [P]. The data $c_T,t_4,\alpha_T$, the sum rules and $\rho>0$ leave $E(1)$ free at $O(1)$, so no
finite-resolution KSS statement follows from controlled data alone (\texttt{code/cft\_etas}). The rate is classical: $E(n)=\int x^n d\nu$ with $x=e^{-2\pi Tt}$ are Hausdorff moments of the relaxation function, the factor
$\prod_{n\le N}n/(n+2\gamma)$ is the M\"untz--Sz\'asz product, and by Smolyak's lemma the bound is the optimal worst-case error of any recovery method
(Micchelli--Rivlin optimal recovery in $H^\infty$; uniqueness from Matsubara data is Baym--Mermin). Extra smoothness or analyticity of $\delta\eta$ at
$\omega=0$ does not help; a one-signed deviation gives exactly $M^{1/(1+\gamma)}\varepsilon^{\gamma/(1+\gamma)}$ (log-convexity of moments), a power of
$\Delta_g$; and a positive relaxation function gives the $M$-free lower bound $4\pi\eta/s\ge E(1)^2/E(2)$ and its truncated-Hausdorff refinements, though
the Einstein relaxation function itself changes sign at late times (\texttt{code/math\_transfer}).

\textbf{Remark 12 (three bounded pieces)} [D]. (a) \emph{Scalar sector.} A canonically normalised scalar with coupling $\gamma\phi C^2$, regular at the
horizon and source-free (either quantisation), has $\phi_h/\gamma=h(\Delta)=\tfrac43\int_0^1\!du\int_0^1\!dt\,[t(1-t)]^{-1}[ut(1-t)/(1-ut)]^{\Delta/3}$,
positive and decreasing, so $h\le C_0=(4\pi/9)\Gamma(\tfrac16)^2/\Gamma(\tfrac13)=16.149$ for $\Delta>\tfrac12$; then
$\delta^{(2)}(4\pi\eta/s)=\gamma^2[24-2m^2h(\Delta)]\in[0,\,6(4+\tfrac34C_0)\gamma^2]$, summed over species. Checks: $24\gamma^2$ for a massless scalar
($\phi=4\gamma/3r^3$), and $\to432\gamma^2/m^2$ at large mass, the integrated-out $I^2$ value [DV, ratio 0.997 at $\Delta=120$]. So an inclusive bound
$\sum_a\gamma_a^2\le K\Delta_g^{-4}$ would control this sector with no mass enhancement, including at the BF point. The whole profile obeys $0<\psi_\Delta\le h(\Delta)$, $\psi_\Delta'\le0$ (maximum principle for $v=r\psi'$), which gives the quadratic entropy
correction $\delta^{(2)}s/s_0=\sum_a q_{\Delta_a}\gamma_a^2$, $q_\Delta=6\int_1^\infty r^{-4}\psi_\Delta\,dr\in(0,2C_0]$ [checks: $q_3=4/3$; $q\to8/m^2$]. In
standard quantisation a nonlinear version holds with constants uniform in masses and multiplicity (Green kernels dominated by the BF kernel, a contraction
in a weighted norm): for $G=\|\boldsymbol\gamma\|_2$ small the shear functional obeys $|\mathscr S-\mathscr S_2|\le\tfrac{15872}{3}vG^3$ with $v$ a bound on
the scalar self-interaction Hessian, so an inclusive $G^2\lesssim\Delta_g^{-4}$ would give $O(\Delta_g^{-4})$ plus $O(\Delta_g^{-6})$. Alternate quantisation
with cubic potentials needs extra boundary data at resonant dimensions. (b) \emph{Coefficient bound.} If a
weak channel $\mathsf S=I+iA/c_T+o(c_T^{-1})$ is contractive, $\operatorname{Im}A\succeq0$, and the matrix typically-real bound gives
$-R^{-2}A_1\preceq A_3\preceq3R^{-2}A_1$ on a disk of radius $R$; a disk $R\ge\kappa\Delta_g^2$ with $\|A_1\|\le M$ uniform would give $\|A_3\|\le3M\kappa^{-2}\Delta_g^{-4}$.
(c) \emph{The logarithm.} A positive impact kernel with Fourier transform $F$ acting on the graviton pole gives $F(0)\log(1/\epsilon)$, and $F(0)>0$
for any nonzero positive kernel, so positivity alone cannot remove the $\log\Delta_g$ of Ref.~\cite{AKT}; a Regge residue holomorphic on
$|\nu|<a\Delta$ with a fourth-order zero at 0 and $|\mathcal R|\le M$ would give $M(3/2a\Delta)^4$ by the Schwarz lemma. The uniform disk, slope and
residue norms are the unproved CFT inputs. Two shortcuts fail as stated: a smeared detector $\int f(\theta)\mathcal E(\theta)$ is positive on all states iff $f\ge0$ (boost a stress-created
state), so kernel differences with equal integrals are never positive; and inverting a positive-imaginary amplitude needs a zero-free disk that a pole
gap does not supply ($-1/z+\lambda/(1-z)$ has a zero at $1/(1+\lambda)$).

\textbf{Remark 13 (real-time data)} [P/D-sketch]. (a) For $t<\beta/2$ the zero-momentum stress commutator is fixed by the convergent Lorentzian thermal OPE
($|t|+|x|<\beta$, causal support), heavy states contribute $\lesssim(2t/\beta)^{\Delta_g}$, and analytic (double-trace) terms give no commutator, so the
thermal coefficient $b_1$ of Remark 2 drops out of the early response only at exact mean-field dimensions: an anomalous dimension
$\delta$ lifts the zero at first order, ${\rm Disc}(-t^2)^{n+\delta/2}=i\pi(-1)^n\delta\,t^{2n}+O(\delta^2)$, and with $\lambda_{ttD}=O(1)$,
$\langle D\rangle_\beta=O(c_T)$, $\delta_D=O(1/c_T)$ this enters at leading connected order, as do light multi-stress primaries (a scalar triple-stress
operator near $\Delta=9$ already has a timelike discontinuity), so the early-time remainder is not $(2t/\beta)^{\Delta_g}$ without control of the
multi-stress data. (b) If the retarded-kernel deviation is
analytic and bounded in a sector $|\arg t|<\alpha$ and $\le\varepsilon$ on $(0,t_2]$, then $|k(t)|\le M^{1-\theta}\varepsilon^\theta$ with
$\theta=(2/\pi)\arcsin((t_2/t)^{\pi/2\alpha})$; with the Einstein quasinormal-mode angles ($\to60^\circ$) this gives $\exp[-2.3(\log1/\varepsilon)^{0.4}]$,
better than any power of the logarithm but not a power of $\Delta_g$ [P]. (c) A one-signed deviation would give a power, but if the first quasinormal mode
that differs from Einstein is not purely damped the deviation changes sign infinitely often; the Einstein relaxation function itself oscillates with
zeros spaced by $\pi/{\rm Re}\,\omega_1$, and even monotone near-horizon deformations give non-one-signed deviations [P given the QNM expansion; DV]. The
remaining CFT-native route is the hyperbolic one: on $\mathbb R\times H^2$ the state at $T_0=1/2\pi R$ is conformal to the vacuum, temperature derivatives
there are integrated vacuum $n$-point functions of $T$, and the planar quantity is the endpoint of a continuation in temperature; all-$n$ vacuum gap bounds
plus uniform analyticity of that continuation would give a power law by Cauchy's estimate (\texttt{code/cft\_etas}).

\textbf{Remark 14 (the hyperbolic continuation)} [D; conditional theorem P]. On $\mathbb R\times H^2_R$ use the strain $\delta g_{ij}=2h\,Y_{ij}$, $Y=dx^2-dy^2$ in
disk coordinates, which is transverse, traceless, $-D^2Y=2Y$, of finite norm $\mathcal N_Y=\pi/6$, and not pure gauge; define
$\eta_Y=-\lim_{\omega\to0}{\rm Im}\,G^R_{O_YO_Y}/(2\mathcal N_Y\omega)$. On the hyperbolic Einstein branch its zero-frequency solution is $H_0=1$, so
$\eta_Y/s=1/4\pi$ at every temperature; at $T_0=1/2\pi R$ the ingoing solution is $(1-i\bar\omega/r)e^{i\bar\omega\,{\rm arccoth}\,r}$ with flux
$-\bar\omega(1+\bar\omega^2)$ [checked symbolically], and since the $T_0$ two-point function is the conformally mapped vacuum one,
$\eta_Y(T_0)=\pi^2C_T/(48R^2)$ is a CFT statement. With the heat capacity $c_H=T\partial_Ts$ (free of the sphere-entropy constant), the Rényi relation
gives $c_H(T_0)=4\pi\eta_Y(T_0)$ exactly, and $8\pi\eta_Y/c_H\to4\pi\eta/s$ at the flat endpoint ($c=2s$ in $d=3$). For Einstein gravity the flat and DC
limits commute uniformly: $\sup_{0\le\lambda\le1}|4\pi\mathcal D_E(\lambda,w)-1|\le Cw^2$ for $|w|<\delta$, $\lambda=r_h^{-2}$ [P]. Temperature derivatives
at $T_0$ are modular cumulants $\kappa_1(K^n;A)$, i.e.\ integrated vacuum correlators with up to $n+2$ stress tensors; the rescaled retarded kernel at
fixed $\omega/T$ also picks up time-rescaling and $q^4$ terms and contact distributions ($\mathcal D[\delta_+']=1$). Conditional theorem: if the
deviations of the normalised $\hat\eta$ and $\hat c$ from Einstein have modular expansions with common summable majorants $\sum b_n\le\epsilon M_\eta$,
$\sum e_n\le\epsilon M_c$, uniform on a small real-frequency region including the DC boundary, then $|4\pi\eta/s-1|\le\epsilon(M_\eta+M_c)/(1-\epsilon M_c)$;
in Einstein gravity the Taylor radius about $T_0$ is $2/\sqrt3>1$, so the flat point lies inside the disk of convergence. Fixed-order bounds do not
suffice ($((1-q)/(1+q))^N$ is uniformly bounded and locally small yet equals 1 at $q=0$), and the absence of a Hawking--Page transition is not an
analyticity theorem. The open input is the all-order bound $\sum_n(|\delta v_n|+|\delta d_n|)\le C\Delta_g^{-4}$ on these integrated multi-point vacuum
energy correlators (\texttt{code/cft\_etas}, THEOREM\_H.md).

\textbf{Remark 15 (the first genuinely five-point coefficient)} [D]. At $T_0$, $\langle\Omega|K_R^mA|\Omega\rangle=\langle\Omega|AK_L^m|\Omega\rangle$ for $A$ in the
right wedge (the boost annihilates the vacuum), so the Taylor coefficients are vacuum correlators across the Rindler horizon with boost-charge light-ray
insertions. On the hyperbolic branch the first-order response of any $F(I,J)$ is explicit and reproduces the planar coefficients at the flat endpoint;
the second coefficient of $\mathcal V$ decomposes as $15\lambda_J+408\lambda_{I^2}-2880\beta_5+O(\lambda^2)$, so its only uncontrolled input is the quintic
coupling. Thermal microcausality sees $\beta_5$ (type-D backgrounds, unlike the type-N shocks of Remark~3): at high momentum the enhanced delay is
$\propto q^2v_1^4K(q)k^4$ with $K=F_{II}-3qF_{IJ}+\tfrac94q^2F_{JJ}$ and a square polarisation factor, reproducing $\alpha_4\ge0$ for $I^2$ and a perfect
square for heavy exchange. For $\beta_5$ it has one sign at every depth and polarisation, so it gives only $\beta_5\ge-O(\Delta_g^{-2})$ (the cone
$b^2\le4ac$ of Remark~4) and leaves the sign that lowers $\eta/s$ free; gradient operators are invisible at this order. Five-point superconvergence,
Tomita--Takesaki positivity and Rényi states also fail (interference terms, three-point reach, no gap dependence). The open input is a two-sided bound
on the inclusive heavy weights $\sum_X|\langle X|T^j\rangle|^2/M_X^2$ for $j\ge3$, beginning with the six-point diagonal.

\subsection{Assumptions of the transport and spectral statements}
\begin{center}
\begin{tabular}{L{6.2cm}L{9.6cm}}
statement & inputs (grade)\\\hline
Shear sum rule & condition (b) of Sec.~\ref{sec:transport}.B [Q; DV]\\
Hall sum rule & H1--H4 [D]; integral form also H5 [D given H5]; exact checks in large-$N$ CS-fermion theory and on the brane\\
Universal disc & both sum rules; collider bound [Q]\\
Theorem SG & A1$'$, A2$'$, inputs of $\tau_*$, condition (b); leading order in $1/c_T$. A2$'$ holds for tree-level Einstein gravity in AdS$_4$ [computed exponents; \texttt{code/a2prime\_einstein}]\\
Fixed-gap and all-gap forms & as for Eq.~(5) of the Letter, resp.\ Eq.~\eqref{eq:fp-allgap} with U1\\
$\mathcal V=18\hat g_3+\dots$; Hall viscosity & bulk effective theory on the brane, tree level [DV]\\
Bound on $|\mathcal V|$ & the above; quartic positivity and L7 for $G_4$ [A]; $|\hat g_4|\le g_4$ [Q, flat space]; Proposition R hypothesis [A]\\
Corrected viscosity bound, disk & the above; $\alpha_4\ge0$, $\tilde\alpha_4^2\le\alpha_4\alpha_4'$ [flat space]; no $\phi W^2$ scalar below the gap\\
EFT-free $\eta/s$ & open [I]\\
\end{tabular}
\end{center}
